\documentclass[11pt,a4paper]{article}
\usepackage[T1]{fontenc}
\usepackage{lmodern,amsmath,amssymb,amsthm,mathtools,booktabs,microtype,xcolor}
\usepackage[margin=26mm]{geometry}
\usepackage[colorlinks=true,linkcolor=blue!45!black,citecolor=blue!45!black,urlcolor=blue!45!black]{hyperref}
\newtheorem{theorem}{Theorem}[section]
\newtheorem{lemma}[theorem]{Lemma}
\newtheorem{proposition}[theorem]{Proposition}
\theoremstyle{remark}\newtheorem{remark}[theorem]{Remark}
\newcommand{\Q}{\mathbb Q}\newcommand{\Z}{\mathbb Z}\newcommand{\HH}{\mathcal H}
\DeclareMathOperator{\Tr}{Tr}
\setlength{\emergencystretch}{2em}
\title{An explicit quartic in 342 variables}
\author{}\date{9 November 2026}
\begin{document}\maketitle
\begin{abstract}
We specify an integer polynomial of degree four in 342 total variables
whose rational input projection is the integers, under the arithmetic
premises of the elliptic source construction. Its complete coefficient
table has 4938 nonzero monomials. Three height detectors use sixteen
local tests expressed by polynomials whose summands have different
valuation classes. Existing formal-logarithm coordinates and nine
additional field squares supply the required powers. The construction retains
the effective abscissa-sum height bound and the explicit fixed parameter
$80989^{54929}$. The polynomial has content one and a largest coefficient
with 1791283 binary digits.
\end{abstract}

\section{Statement and coefficient specification}
The rational input projection of a polynomial $P\in\Z[X,\boldsymbol Y]$
is the set of $x\in\Q$ for which $P(x,\boldsymbol y)=0$ has a solution
in rational values of $\boldsymbol Y$. Variable counts include the input
$X$, and degree always means total degree. The coefficient height $H(P)$
is the largest absolute integer coefficient.

For an integer $N\ge1$ put $A_N=80989^{2N+1}$. The exact tables
\nolinkurl{quadratic-family342.json} and
\nolinkurl{polynomial-family342.json} specify 314 primitive integer
quadratic rows $Q_{N,j}$ in 342 coordinates and
\[
                    P_N=\sum_{j=1}^{314}Q_{N,j}^2.
\]
The powers of $A_N$ are recorded separately in the tables solely to
abbreviate integer coefficients. The parameter is not a variable.

\begin{theorem}\label{thm:main}
Under the arithmetic premises in Section~\ref{sec:premises}, for every
integer $N\ge27464$ and every $x\in\Q$,
\[
 x\in\Z\quad\Longleftrightarrow\quad
 \exists\boldsymbol y\in\Q^{341},\ P_N(x,\boldsymbol y)=0.
\]
Every such polynomial has degree exactly four, content one, and 4938
nonzero coordinate monomials. Its coefficient height satisfies
\[
                         H(P_N)=80989^{4N+2}+1.
\]
In particular, substituting $A=80989^{54929}$ specifies a single
polynomial $P_{27464}$ with the asserted projection. Its largest
coefficient has exactly 1791283 binary digits. Algebraically, the support
count holds for every $N\ge1$, and the height formula already holds
for every $N\ge23$.
\end{theorem}

\section{Arithmetic premises}\label{sec:premises}
Put
\[
 K=\Q(\rho),\quad\rho^2=2,\quad d=3-\rho,\quad
 E:y^2=x^3-d^2x,
\]
\[
 R_0=(1+\rho,2\rho),\qquad
 P=2R_0=\left(\frac{57-28\rho}{8},\frac{-660+417\rho}{32}\right).
\]
We retain the following arithmetic contract, stated with its exact
source locations in \cite[Section~2]{Source}.
\begin{enumerate}
\item The group and index-height assertions are
\[
 2E(K)=\Z P,\qquad h(x([n]P))\ge n^2-6\quad(n\ne0).
\]
Here $h$ is absolute logarithmic Weil height, and $\sigma$ denotes
conjugation on $K/\Q$.
\item For $t=u/v$ in lowest terms, $v>0$, outside
$\{0,1,2401,-343/32\}$, let $M_5(t)$ be the product of primes
$p>76176$ with a positive odd valuation on at least one of
\[
             v,\quad u,\quad u-v,\quad u-2401v,\quad32u+343v.
\]
The contact-height input is
\[
 h(t)\le\HH M_5(t)^{19},\qquad
 \HH=2^{110}210^{1800}\left(\prod_{p\le76176}p\right)^{19}+128
       <2^{2093448}.
\]
\item The source pole equations exclude negative odd valuations of their
nonzero input $b$ at primes outside $\{2,3,17\}$. They admit the specified
specialization when $0<b<17$, its reduced denominator is a square,
$b$ is a square in $\Q_2,\Q_3,\Q_{17}$, and its valuations there are
at least $3,1,2$, respectively. Their exact normalized seven-coordinate
form is retained in the coefficient table.
\item The simultaneous prime-pattern and conjugate reduction-order input
holds with arbitrary fixed odd avoidance modulus, the prescribed real
approximation, and the finite congruences simultaneously. The selected
primes are $1$ modulo $8\cdot3\cdot17$, with
$v_2(d_-)=6$, $v_2(d_+)=2$ and
$\gcd(d_-(p),d_+(p'))=4$. Its quantitative weighted form is retained:
for fixed distinct coefficients $c_j$, admissible patterns
$p_j=u+c_jw^2$, $N'/2\le u\le N'$,
$1\le w\le\tau(N')\sqrt{N'}$, have total weight at least
$c_0(N')^{3/2}\tau(N')$, individual weights
$O((\log N')^k)$, and $\tau(N')$ may decrease arbitrarily slowly
after the fixed avoidance and congruence conditions are imposed.
\end{enumerate}
These are arithmetic premises, not conclusions of coefficient
verification. The weighted selection and
separation used here follow from the cited analytic estimates as proved
in \cite{Selection}; no additional prime-distribution conjecture is
introduced by the local weight replacement. The effective sum-height theorem used below is proved
in \cite{Sums}; it is not an extra conjectural premise.

\section{The local tests and the quadratic system}
The four root charts are
\[
 x=h^2,\qquad j^2=x^2-d^2,
\]
and give nonzero integer multiples $Q_i$ of $P$. The squareclass
argument establishing this assertion, including its converse and all
torsion exclusions, is proved in \cite[Section~2]{Baseline}.
Write their labels as $n,a,t,b$ and set
\[
 B=\max(|n|,|a|,|t|,|b|),\qquad
 S=\sum_{i=1}^4x(Q_i)=S_0+S_1\rho.
\]
The three positive rational square scales are
\[
                         1,\qquad S_0^2,\qquad(S_0+S_1)^2.
\]
They are nonzero because $S_0>12$ and $S_0>\sqrt2|S_1|$.

For each detector the contact input is $t_*=a_*\lambda$ for one of
these scales. Put $\delta=t_*-1$,
\[
 F_*(t_*)=t_*(t_*+1/2)(t_*+1)(t_*-1/2399)(t_*+32/407),
\]
\[
 R_*=\frac{F_*(t_*)}{a_*^5\delta^6},\qquad
 q=R_*^{-1}\delta^{-4},\qquad q_m=q^{2m-1}.
\]
These rational formulas are used only when their denominators are
nonzero; the actual system consists of the polynomial rows specified
in the table. The exceptional value $t_*=1$ has empty contact radical
and is explicitly covered by the local proof in \cite{Packing}.

Write $f(Q)=1/x(Q)$, extended by zero at the identity. The auxiliary
points satisfy $U_0=H$, $V_0=H-P$, and $Q_i=U_i-V_i$ as group
differences on $E$. The exact addition rows and identity-inclusive
charts of \cite{Baseline} enforce these relations.

Define $A_i=f(U_i)$, $B_i=\sigma f(V_i)$, and
\[
 \begin{aligned}
 C_i&=L_i-L'_i\quad(1\le i\le4),& I&=L_1-\beta,\\
 D_H&=f(H),&G&=L_3^2-L_2(L_2-2)L_3L_4+(L_2-2)^2L_4^2-1.
 \end{aligned}
\]
The alternate ratios $L'_i$ already include conjugation. The nine tests
$A_i,B_i$ have order one, $C_1,I$ have order ten, and $D_H,C_2,C_3,C_4,G$
have order forty. For an ordered list of tests, use the factor
\[
 \mathcal T(D_0,\ldots,D_{s-1})
 =1+\sum_{j=0}^{s-1}q_{m(D_j)}^{2^j}D_j^{2^{j+1}},
\]
where $m(D_j)$ is its stated order. Thus different required orders can
occur within one factor.

The detector uses the product of the factors for these five lists:
\[
 \begin{gathered}
 (D_H,C_1,B_1,B_0),\quad(C_2,I,B_2),\quad(C_3,A_2,B_3),\\
 (C_4,A_3,B_4),\quad(G,A_4,A_1).
 \end{gathered}
\]
Each test occurs exactly once. The additional field powers are
\[
 C_1^4,I^4,A_2^4,A_3^4,A_4^4,\quad B_1^8,B_0^{16},\quad A_1^4,A_1^8.
\]
In particular the sixteenth power of $B_0$ is obtained by squaring its
existing eighth power, and the eighth power of $B_1$ by squaring its
existing fourth power. All additional field powers use quadratic rows.
The full normalization formulas and exact coordinate accounting are
in \cite{Packing}.

At a contact prime, a nonzero term in position $j$ has valuation
$2^j$ times an odd integer. These valuations are pairwise distinct
and different from zero. The unique-minimum argument and the literal
trace-clearing rows give every assigned test order at the same place;
this is \cite[Proposition~2.3]{Shared}. The argument also proves the
nonvanishing needed for the trace quotient. The rational logarithms,
point charts, and scalar equations retain their stated forms.


\subsection{A quadratic presentation of the local weights}
For each detector, write $F_*(t_*)=\prod_{\theta\in\Theta}(t_*-\theta)$
and put $q=R_*^{-1}\delta^{-4}$. The weights used here
are $q_m=q^{2m-1}$, so $q_{10}=q^{19}$ and $q_{40}=q^{79}$.
The quadratic rows are
\[
 \begin{gathered}
 q_2=q^2,\quad q_4=q_2^2,\quad c q_4=q,\quad q_8=q_4^2,\\
 c q_{19}=q_8^2,\quad q_{38}=q_{19}^2,\quad c q_{79}=q_{38}^2.
 \end{gathered}
\]
At every consumed contact prime, $\delta$ is a unit and $q\ne0$.
The rows then force $c=q^{-3}$ and each indicated power uniquely.
No inverse coordinate other than $c$ is needed. If $q=0$, the retained
relations $\delta^3q=e$ and $F_*(t_*)e=(a_*\delta)^5$ force
$\delta=0$; the contact radical is empty in that case.
This covers the potentially degenerate stratum without using a rational
formula outside its domain. The full local extraction proof from the
literal trace equations is \cite[Section 2]{Shared}.

At contacts these powers have exactly the required valuations:
$v(q_m)=(1-2m)v_p(R_*)$. The ordered factors therefore supply the same
orders $1,10,40$. At prepared finite primes their simultaneous smallness
is established directly in \cite[Section 3]{Shared}, by enlarging a fixed
denominator before prime selection. Thus the replacement changes neither
the formal logarithm comparisons nor their required depths.


\subsection{Primitive normalization of the scalar contact product}
Let $z_j$ denote a coordinate by its source identifier. The actual
contact product is the affine expression
\[
 \begin{aligned}
 r={}&-103171881284-2103220824z_{248}-21050600z_{487}\\
     &-7629700z_{488}-45771136z_{489}.
 \end{aligned}
\]
We use $r$ itself as the coordinate in place of $z_{489}$. This is an
invertible rational affine change: its coefficient of $z_{489}$ is
nonzero, and every other coordinate, including the input, is fixed.
The source coordinate occurs only in its own defining row and the
scalar normalizer row. After substitution and primitive normalization,
the latter is exactly
\[
                  R_s+80989^2rR_s-A_Nr=0.
\]
Every equation remains quadratic. Reverse substitution reproduces both
source rows exactly up to a nonzero rational scalar, so their rational
solution sets correspond bijectively. Taking sums of squares preserves
the required conjunction over $\Q$.

In the resulting quartic the coefficient of $r^2$ is $A_N^2+1$.
Every other coefficient is either a nonzero constant or a nonzero constant
multiple of $A_N$. Thus no coefficient vanishes for a positive parameter,
and the support counts hold for every $N\ge1$.
Let $C_0$ be the largest absolute constant coefficient. Its exact integer
value is recorded in the certificate and has 1520 binary digits. The
largest absolute coefficient multiplying $A_N$ is
\[
                       L_0=2\cdot80989^2=13118436242.
\]
The complete integer comparison verifies
\[
 80989^{94}+1>C_0,\qquad80989^{47}>L_0,
 \qquad80989^{90}+1<C_0.
\]
Consequently, for both families,
\[
                        H(P_N)=80989^{4N+2}+1
                        \qquad(N\ge23).
\]
Indeed this coefficient exceeds both $C_0$ and $L_0A_N$ at $N=23$,
and continues to do so as $N$ increases. Since $C_0>L_0^2$, the exact
height for any positive parameter is $\max(C_0,A_N^2+1)$.
This is an algebraic height threshold; the sufficient parameter in the
342-variable defining theorem remains $27464$, and the corresponding
threshold for the 297-variable family remains ineffective.

\section{Soundness with the explicit parameter}
The input label and the scalar normalizer are
\[
 \begin{gathered}
 \beta=51X^2+2504,\qquad
 r=\beta(\beta-1)(\beta-2401)(32\beta+343)(2\beta-1),\\
 R_s=\frac{A_Nr}{1+80989^2r},\qquad c_s=R_s+9792.
 \end{gathered}
\]
The retained scalar argument gives, with $c=|a-1|$,
\[
 c\ge80989^{L_N},\qquad
 L_N=\left\lfloor\frac{187(2N-1)}8\right\rfloor+1,\qquad
 h(\beta)\le\HH c^{19/24}.
\]
Its exceptional-point exclusions precede every use of the height
detectors and are unaffected by their polynomial factors. The integer
certificates in \cite{Baseline} give, for $N\ge27464$,
\begin{equation}\label{eq:margins}
 |a|>(40\HH)^{10},\qquad
 \log(48\HH)<|a|/40,\qquad h(\beta)<|a|/40.
\end{equation}

The effective sum-height theorem yields
\[
 B^2\le2h(S_0)+2h(S_0+S_1)+82.
\]
Combining it with the contact-height input for the three scales gives:
if a nonzero integer $J$ is divisible by the $k$th powers of all three
contact radicals, then
\begin{equation}\label{eq:height}
                          B^2\le5\HH |J|^{19/k}.
\end{equation}
The formal logarithm comparison uses only the orders $1,10,40$,
which the displayed factors supply unchanged. Therefore the fortieth
powers of the contact radicals divide
\[
 \Delta=t^2-a(a-2)tb+(a-2)^2b^2-1.
\]
If $n-\beta\ne0$, their tenth powers divide its nonzero reduced
numerator $J$.

\begin{proof}[Proof of soundness in Theorem~\ref{thm:main}]
A rational zero of $P_N$ makes all its quadratic rows zero.
If $\Delta\ne0$, then $|\Delta|\le14B^4$, and \eqref{eq:height}
with $k=40$ gives
\[
 B^2\le5\HH(14B^4)^{19/40}<20\HH B^{19/10}.
\]
Thus $B^{1/10}<20\HH$, contradicting $B\ge|a|$ and
\eqref{eq:margins}. Hence $\Delta=0$.

The integral Lucas descent in \cite[Section~5]{Source} applies to all
nonzero integer $b$ and $|a|>7$, giving $\log B>|a|-3$.
If $n-\beta\ne0$, its reduced numerator satisfies
$|J|\le2B\exp(h(\beta))$. Equation~\eqref{eq:height} with $k=10$
would imply
\[
 \log B<10\log(20\HH)+19h(\beta)<29|a|/40,
\]
contrary to the Lucas bound. Thus $n=\beta\in\Z$. Write $X=u/v$
in lowest terms with $v>0$. Integrality of $51X^2+2504$ implies
$v^2\mid51$. Since $51$ is squarefree, $v=1$ and $X\in\Z$.
\end{proof}

\section{Completeness and exact verification}
Fix an integer input and the scalar parameter. The retained scalar and
Lucas construction supplies the integer labels and nonzero root points.
Choose compact smaller real and finite-place neighborhoods avoiding the
poles of the retained rational tests. The tests then have fixed valuation
lower bounds at every prime in the finite preparation set $S$.
All scales and neighborhoods are fixed before the contact primes.

Choose a positive integer $D$ that clears the contact coefficients
$D^2\theta/a_*$ and has sufficiently large valuations at the primes in
$S$. In the simultaneous prime pattern put
\[
 \lambda=\frac{u}{D^2w^2}.
\]
The zero contact makes $u$ a selected prime, outside $S$. Therefore
$v_p(\lambda)\le-2v_p(D)$. Making this upper bound sufficiently negative
gives $v_p(t_*)<\min(0,v_p(\theta):\theta\ne0)$ and hence
\[
 v_p(F_*(t_*))=5v_p(t_*),\quad
 v_p(e)=5v_p(a_*),\quad
 v_p(q)=2v_p(a_*)-3v_p(\lambda).
\]
The latter is uniformly arbitrarily large. Thus all new weighted test
terms are arbitrarily small at the prepared finite places. Enlarging $D$
is a fixed dilation of the prime pattern and is allowed by the retained
selection theorem; see \cite{Selection}.

At each selected contact prime, the designated primary place gives
all test orders. The distinct-valuation argument at the opposite place
makes $\Tr(W^{-1})$ a unit congruent to $1$ or $2$ modulo $p>2$.
Consequently the pole input $R_*\Tr(W^{-1})^{-2}$ has valuation one.
The exact product formula
\[
 R_*=\frac{\prod_{\theta\in\Theta}(u-(\theta/a_*)D^2w^2)}
               {(D^2w^2)^5\delta^6}
\]
shows that its valuation is even at every other prime. This includes
all primes dividing $\delta$; no unit assertion about $\delta$ away from
contacts is required. Its reduced denominator is therefore a square.
The selected-prime congruences give the required local squareclasses
at $2,3,17$, and the fixed choice of $D$ gives their required valuations.

Uniformly on the fixed compact real neighborhoods,
\[
 R_*\sim a_*^{-6}\lambda^{-1},\qquad q\sim a_*^2\lambda^{-3},
 \qquad \Tr(W^{-1})\longrightarrow2.
\]
All factors tend to $1$, and the pole input tends to $0^+$. It therefore
lies in $(0,17)$ for sufficiently large $\lambda$. The simultaneous
selection supplies such parameters after all fixed data have been chosen.
The pole-realization premise gives fresh witnesses, and the rational
weight identities fill every new quadratic row. Equal scales may reuse
witnesses. This is the full simultaneous completeness argument of
\cite[Proposition 3.2]{Shared}; it imposes no new lower bound on the
scalar parameter $N$.

The exact reproduction programs are
\begin{verbatim}
python build_quartic_shared_weights.py
python check_quartic_shared_weights.py
\end{verbatim}
The checker imports no producer. It checks every changed and retained
row, all coordinate incidences, and every coefficient of the quartic by
an independent unordered-pair expansion. It verifies degree four,
content one and the coefficient-height bound. These finite checks do
not prove the retained global arithmetic inputs.

The exact table contains 314 rows, 342 total coordinates, 4939 terms
before parameter powers are collected, and 4938 ordinary monomials for
every $N\ge1$. The coordinate $q_8$ also supplies the factor term in position three;
the same value is used in the shared weight presentation.
Direct integer evaluation at $A=80989^{54929}$ verifies that the largest
coefficient has 1791283 binary digits. It is the coefficient of $r^2$.

\begin{thebibliography}{9}
\bibitem{Shared}
\emph{Common local weights and quartics in 297 and 342 variables},
9 November 2026, Sections 2--4: local extraction, simultaneous preparation,
nonzero domains and exact quadratic weight presentation.
\href{../local_tests/shared_local_weights_20261109_v12.pdf}{PDF};
source: \nolinkurl{shared_local_weights_20261109_v12.tex}.
\bibitem{Selection}
\emph{Separation and denominator conditions in quadratic prime patterns},
9 November 2026, weighted separation from Sun's stated analytic estimates.
\href{../prime_selection/prime_selection_20261109_v12.pdf}{PDF};
source: \nolinkurl{prime_selection_20261109_v12.tex}.

\bibitem{Packing}
\emph{Simultaneous valuation conditions and explicit quartics},
9 November 2026, accompanying proof and independent exact verification.
\href{../local_tests/mixed_valuation_tests_20261109_v11.pdf}{PDF}; source: \nolinkurl{mixed_valuation_tests_20261109_v11.tex}.
\bibitem{Baseline}
\emph{A quartic in 370 variables with specified integer coefficients},
8 November 2026, complete quadratic specification and effective height proof.
\href{quartic370_20261108_v10.pdf}{PDF}; source: \nolinkurl{quartic370_20261108_v10.tex}.
\bibitem{Source}
\emph{A forced prime condition and an explicit quartic in 392 variables},
8 November 2026, arithmetic input contract and scalar/Lucas arguments.
\href{quartic392_20261108_v8.pdf}{PDF}; source: \nolinkurl{quartic392_20261108_v8.tex}.
\bibitem{Sums}
\emph{Effective heights of sums of elliptic abscissas}, 8 November 2026,
effective four-term sum-height theorem and its application to two rational
linear forms.
\href{../effective_sums/effective_elliptic_sums_20261108_v10.pdf}{PDF}; source: \nolinkurl{effective_elliptic_sums_20261108_v10.tex}.
\bibitem{Glazer}
E.~Glazer, \emph{An explicit quartic defining the integers in the rationals},
Version~13, manuscript and coefficient data, commit
\nolinkurl{5bbd74f6a479776a332cbe036809d5dddae63720},
\url{https://github.com/elliotglazer/diophantine-z-q}.
\end{thebibliography}
\end{document}
